\begin{dang}{Năng lượng trong phóng xạ}
\Anlythuyet{
\begin{itemize}
	\item Vận dụng bảo toàn năng lượng và bảo toàn động lượng trong phản ứng hạt nhân
	\item Cho phân rã: \begin{align*} \Hn{{Z_1}}{{A_1}}{X}\rightarrow \Hn{{Z_2}}{{A_2}}{Y}+ \text{(tia phóng xạ)}\ \Hn{{Z_3}}{{A_3}}{Z}\ \ \end{align*}
	\lefthand\ Do trước khi phân rã hạt nhân đứng yên nên theo định luật bảo toàn động lượng ta có:
	\begin{align*}
	\vec {p}_Y+\vec {p}_Z=\vec 0\Rightarrow p^2_Y=p^2_Z \Rightarrow 2m_YK_Y=2m_ZK_Z \Rightarrow \boder{\dfrac{K_Y}{K_Z}= \dfrac{v_Y}{v_Z}=\dfrac{m_Z}{m_Y}}
	\end{align*}
\end{itemize}	
}
\end{dang}